\documentclass[tikz,border=4pt]{standalone}
\usepackage{ctex}
\usepackage{amsmath}
\usetikzlibrary{calc, arrows.meta, angles, quotes}

\begin{document}
\begin{tikzpicture}[
  x={(1.8cm,0cm)},
  y={(0cm,1.8cm)},
  z={(-0.52cm,-0.42cm)},
  line cap=round,
  vec/.style={-{Stealth[length=7pt, width=4pt]}, thick},
  vecblue/.style={-{Stealth[length=7pt, width=4pt]}, thick, blue},
]

% 空间两向量夹角示意
% 共起点 O，两个向量 a 和 b，夹角 θ

\coordinate (O) at (0,0,0);

% 向量 a：沿 x 轴方向
\coordinate (Aend) at (1.4,0,0);

% 向量 b：偏转一定角度（在 xy 平面内取 cos60=0.5 方向）
% 实际夹角约 50°，视觉清晰
\coordinate (Bend) at (0.8,1.0,0);

% 画向量
\draw[vec, black] (O) -- (Aend);
\node[right] at (Aend) {$\vec{a}$};

\draw[vecblue] (O) -- (Bend);
\node[above] at (Bend) {$\vec{b}$};

% 夹角弧（从 Aend 方向到 Bend 方向）
% 用 2D 投影角度
% a 方向角 ≈ 0°，b 方向角 ≈ atan2(1.0,0.8) ≈ 51.3°
\draw[black, thin] (O) ++(0.35,0,0) arc[start angle=0, end angle=51.3, radius=0.35cm];
\node[right, font=\small] at (0.38,0.17,0) {$\theta$};

% 原点标注
\node[below left] at (O) {$O$};

% 辅助：虚线连接两向量终点（展示夹角位置）
\draw[gray, dashed] (Aend) -- (Bend);

% 模长标注
\node[below, font=\small] at ($(O)!0.55!(Aend)$) {$|\vec{a}|$};
\node[left, font=\small] at ($(O)!0.55!(Bend)$) {$|\vec{b}|$};

% 说明文字
\node[below] at (0.7,-0.28,0) {$\vec{a}\cdot\vec{b}=|\vec{a}||\vec{b}|\cos\theta$\quad（$0\le\theta\le\pi$）};

\end{tikzpicture}
\end{document}
